% Partie : méthodes
% Chapitre : démonstrations
% Section : démonstration par l'absurde

\documentclass[../../../include/open-logic-section]{subfiles}

\begin{document}
\olfileid[fr]{mth}{prf}{con}

\olsection{Démonstration par l'absurde}

La démonstration par l'absurde est d'abord un schéma
d'inférence destiné à établir des énoncés négatifs.
Supposons que vous vouliez montrer qu'un énoncé~$p$
est \emph{faux}, autrement dit établir~$\lnot p$.
La stratégie la plus prometteuse est de
(a) supposer $p$ vrai, puis de (b) montrer
que cette hypothèse conduit à quelque chose
que vous savez être faux. Ce résultat peut
contredire $p$ lui-même ou une autre hypothèse
de l'énoncé étudié. Par exemple, pour démontrer
«~si $q$, alors $\lnot p$~», on suppose $q$
et on en déduit~$\lnot p$. Si l'on établit
$\lnot p$ par l'absurde, on ajoute
temporairement l'hypothèse~$p$. Si, avec cette
hypothèse, on déduit $\lnot q$, cette conclusion
contredit l'hypothèse~$q$. Il faudra bien sûr aussi utiliser
d'autres schémas d'inférence et déplier
des définitions pour parvenir à la
contradiction. Voici un exemple.

\begin{prop}
  Si $A \subseteq B$ et $B = \emptyset$,
  alors $A$ n'a aucun !!{element}.
\end{prop}

\begin{proof}
  Supposons $A \subseteq B$ et
  $B = \emptyset$. Nous voulons montrer
  que $A$ n'a aucun !!{element}.
  \begin{quote}
    L'énoncé est une implication :
    nous supposons son antécédent et
    cherchons à démontrer son conséquent.
    Celui-ci dit que $A$ n'a aucun
    !!{element}. Reformulons-le plus
    explicitement : il est faux qu'il
    existe un~$x \in A$.
  \end{quote}
  $A$ n'a aucun !!{element} si et
  seulement s'il n'existe pas de~$x$
  tel que $x \in A$.
  \begin{quote}
    Nous voulons donc démontrer un
    énoncé négatif~$\lnot p$ :
    il est faux qu'il existe un
    $x \in A$. Pour procéder par
    l'absurde, supposons l'énoncé
    positif correspondant~$p$,
    c'est-à-dire l'existence d'un
    $x \in A$, puis tirons-en
    une contradiction. Nous annonçons
    cette démarche par «~supposons,
    pour obtenir une contradiction,
    que…~», ou simplement par
    «~supposons le contraire~».
  \end{quote}
  Supposons le contraire :
  il existe un $x \in A$.
  \begin{quote}
    C'est notre nouvelle hypothèse
    pour obtenir une contradiction.
    Nous avons aussi $A \subseteq B$
    et $B = \emptyset$. La première
    donne $x \in B$ :
  \end{quote}
  Puisque $A \subseteq B$, $x \in B$.
  \begin{quote}
    Mais comme $B = \emptyset$,
    tout !!{element} de~$B$,
    notamment $x$, doit aussi être
    !!a{element} de~$\emptyset$.
  \end{quote}
  Puisque $B = \emptyset$,
  $x \in \emptyset$. C'est une
  contradiction : par définition,
  $\emptyset$ n'a aucun !!{element}.
  \begin{quote}
    La démonstration est déjà achevée :
    nos hypothèses ont conduit à une
    contradiction et ne peuvent donc
    être toutes vraies. Puisque les
    deux premières, $A \subseteq B$
    et $B = \emptyset$, ne sont pas
    mises en cause, c'est la dernière,
    l'existence d'un $x \in A$, qui
    doit être fausse. Nous pouvons
    expliciter cette conclusion.
  \end{quote}
  Ainsi, l'hypothèse qu'il existe
  un $x \in A$ est fausse ; par
  l'absurde, $A$ n'a donc aucun
  !!{element}.
\end{proof}

En logique classique, tout énoncé positif
est équivalent à un énoncé négatif :
$p$ si et seulement si $\lnot\lnot p$.
On peut donc aussi établir un énoncé
positif \emph{indirectement}, par l'absurde.
Pour démontrer~$p$, considérons l'énoncé
négatif $\lnot\lnot p$. Si l'hypothèse
$\lnot p$ conduit à une contradiction,
nous avons démontré $\lnot\lnot p$,
et par conséquent~$p$.

Dans l'exemple précédent, notre objectif
était négatif : montrer que $A$ n'a
aucun !!{element}. L'hypothèse ajoutée
pour la démonstration par l'absurde,
l'existence d'un $x \in A$, était
donc positive. Elle nous a donné
un point de départ concret,
l'appartenance hypothétique $x \in A$,
à partir de laquelle nous avons
raisonné jusqu'à $x \in \emptyset$.

Pour démontrer indirectement un énoncé
positif, on commence au contraire par
une hypothèse négative. Mais on peut
souvent reformuler facilement un énoncé
positif comme un énoncé négatif, et
inversement. La démonstration précédente
aurait été essentiellement la même si
nous avions cherché à établir
«~$A = \emptyset$~» plutôt que
«~$A$ n'a aucun !!{element}~».
Par définition de~$=$,
«~$A = \emptyset$~» est un énoncé
universel : il se déplie en
«~tout !!{element} de~$A$ est
!!{element} de~$\emptyset$,
et réciproquement~».
Mais il est facile de voir qu'il
équivaut à l'énoncé négatif
«~il n'existe pas de $x \in A$~».

Il est parfois plus facile de partir
de l'hypothèse $\lnot p$ que de
démontrer directement~$p$. Même
lorsqu'une démonstration directe serait
aussi simple, voire plus simple,
certains préfèrent l'argument indirect.
Si la double négation vous déroute,
retenez les deux points de départ.
Pour démontrer $\lnot p$, on suppose
$p$ et on en tire une contradiction.
Pour démontrer $p$ en logique classique,
on tire une contradiction de sa
\emph{négation}~$\lnot p$.

\begin{prop}
$A \subseteq A \cup B$.
\end{prop}

\begin{proof}
  Nous voulons montrer
  $A \subseteq A \cup B$.
  \begin{quote}
    À première vue, c'est un
    énoncé positif : tout
    $x \in A$ appartient aussi
    à $A \cup B$. Sa négation
    dit qu'un certain $x \in A$
    n'appartient pas à $A \cup B$.
    Nous pouvons donc supposer
    cette négation et montrer
    qu'elle mène à une contradiction.
  \end{quote}
  Supposons le contraire,
  c'est-à-dire
  $A \nsubseteq A \cup B$.
  \begin{quote}
    Par définition,
    $A \subseteq A \cup B$
    signifie que tout $x \in A$
    appartient aussi à $A \cup B$.
    Pour comprendre
    $A \nsubseteq A \cup B$,
    nions correctement cet
    énoncé universel :
    il est faux que tout
    $x \in A$ appartienne
    à $A \cup B$ si et
    seulement s'il existe
    \emph{un certain}~$x \in A$
    tel que $x \notin A \cup B$.
    Il faut être attentif :
    beaucoup de tentatives de
    démonstration par l'absurde
    échouent parce qu'elles
    nient incorrectement un
    énoncé comme «~tous les
    $A$ sont des $B$~».
    Autrement dit,
    $A \nsubseteq A \cup B$
    si et seulement s'il
    existe un $x$ tel que
    $x \in A$ et
    $x \notin A \cup B$.
    La suite est facile.
  \end{quote}
  Il existe donc un $x \in A$
  tel que $x \notin A \cup B$.
  Par définition de~$\cup$,
  $x \in A \cup B$ si et
  seulement si $x \in A$
  ou $x \in B$. Puisque
  $x \in A$, nous avons
  $x \in A \cup B$.
  Cela contredit
  $x \notin A \cup B$.
\end{proof}

\begin{prob}
Démontrez \emph{indirectement}
que $A \cap B \subseteq A$.
\end{prob}

\begin{prop}
Si $A \subseteq B$ et
$B \subseteq C$, alors
$A \subseteq C$.
\end{prop}

\begin{proof}
  Supposons $A \subseteq B$
  et $B \subseteq C$. Nous
  voulons montrer
  $A \subseteq C$.
  \begin{quote}
    Procédons indirectement :
    supposons la négation
    de notre objectif.
  \end{quote}
  Supposons donc le contraire,
  $A \nsubseteq C$.
  \begin{quote}
    Comme auparavant,
    $A \nsubseteq C$ signifie
    qu'il n'est pas vrai
    que tout $x \in A$
    appartienne aussi à~$C$ :
    il existe donc un
    $x \in A$ tel que
    $x \notin C$. Avec
    l'habitude, ce dépliage
    des négations deviendra
    immédiat.
  \end{quote}
  Autrement dit, il existe
  un~$x$ tel que
  $x \in A$ et
  $x \notin C$.
  \begin{quote}
    Nous pouvons maintenant
    tirer une contradiction,
    en utilisant bien sûr
    les deux autres hypothèses.
  \end{quote}
  Puisque $A \subseteq B$,
  $x \in B$. Puisque
  $B \subseteq C$,
  $x \in C$. Cela
  contredit $x \notin C$.
\end{proof}

\begin{prop}
Si $A \cup B = A \cap B$,
alors $A = B$.
\end{prop}

\begin{proof}
  Supposons
  $A \cup B = A \cap B$.
  Nous voulons montrer
  $A = B$.
  \begin{quote}
    Le début de la
    démonstration est
    désormais familier :
  \end{quote}
  Supposons, pour obtenir
  une contradiction,
  que $A \neq B$.
  \begin{quote}
    Notre hypothèse
    provisoire est
    $A \neq B$. Or
    $A = B$ si et
    seulement si
    $A \subseteq B$
    \emph{et}
    $B \subseteq A$.
    Par négation,
    $A \neq B$ si et
    seulement si
    $A \nsubseteq B$
    \emph{ou}
    $B \nsubseteq A$.
    Il faut être
    attentif lorsqu'on
    nie une conjonction !{}
    Pour tirer une
    contradiction de
    cette disjonction,
    distinguons deux
    cas et obtenons-en
    une dans chacun.
  \end{quote}
  $A \neq B$ si et
  seulement si
  $A \nsubseteq B$
  ou $B \nsubseteq A$.
  Distinguons deux cas.
  \begin{quote}
    Dans le premier,
    $A \nsubseteq B$ :
    un certain $x$
    appartient à~$A$,
    mais pas à~$B$.
    L'intersection
    $A \cap B$ contient
    les !!{element}s
    communs à $A$
    et à $B$ :
    ce qui manque
    dans l'un ne
    peut appartenir
    à leur intersection.
    L'union $A \cup B$
    contient, elle,
    tout ce qui
    appartient à
    l'un ou l'autre.
    Ainsi,
    $x \in A \cup B$
    mais
    $x \notin A \cap B$,
    d'où
    $A \cap B \neq A \cup B$.
  \end{quote}

  Cas 1 :
  $A \nsubseteq B$.
  Il existe alors
  un $x \in A$
  tel que
  $x \notin B$.
  Puisque
  $x \notin B$,
  $x \notin A \cap B$.
  Puisque $x \in A$,
  $x \in A \cup B$.
  Donc
  $A \cap B \neq A \cup B$,
  contrairement à
  l'hypothèse
  $A \cap B = A \cup B$.

  Cas 2 :
  $B \nsubseteq A$.
  Il existe alors
  un $y \in B$
  tel que
  $y \notin A$.
  Comme précédemment,
  $y \in A \cup B$,
  mais
  $y \notin A \cap B$.
  Donc
  $A \cap B \neq A \cup B$,
  contrairement à
  l'hypothèse
  $A \cap B = A \cup B$.
\end{proof}

\end{document}
